\documentclass{article}

\usepackage{amsmath,amssymb}
\usepackage{xurl}
\usepackage[hidelinks]{hyperref}
\setlength{\emergencystretch}{2em}

% Shared commands for notes in docs/paper-gaps/.

\DeclareUnicodeCharacter{2080}{\ifmmode{}_{0}\else\textsubscript{0}\fi}
\DeclareUnicodeCharacter{2081}{\ifmmode{}_{1}\else\textsubscript{1}\fi}
\DeclareUnicodeCharacter{2082}{\ifmmode{}_{2}\else\textsubscript{2}\fi}
\DeclareUnicodeCharacter{2099}{\ifmmode{}_{n}\else\textsubscript{n}\fi}

\newcommand{\Lean}{\textsc{Lean}}
\ExplSyntaxOn
\NewDocumentCommand{\leanid}{m}{
  \ifmmode
    \textsf{\detokenize{#1}}
  \else
    \group_begin:
    \urlstyle{sf}
    \tl_set:Nn \l_tmpa_tl {#1}
    \tl_replace_all:Nnn \l_tmpa_tl {\_} {_}
    \exp_args:NV \path \l_tmpa_tl
    \group_end:
  \fi
}
\ExplSyntaxOff
\providecommand{\tr}{\operatorname{tr}}
\newcommand{\tnleanIssueBase}{https://github.com/LionSR/TNLean/issues/}
\newcommand{\tnleanPrBase}{https://github.com/LionSR/TNLean/pull/}
\newcommand{\tnleanDocBase}{https://sirui-lu.com/TNLean/docs/}
\newcommand{\ghissue}[1]{\href{\tnleanIssueBase#1}{GitHub issue \##1}}
\newcommand{\ghpr}[1]{\href{\tnleanPrBase#1}{PR \##1}}
\newcommand{\doclink}[1]{\href{\tnleanDocBase#1.html}{\path{#1}}}

% Dirac notation and the identity operator, as in blueprint/src/macros/common.tex.
% \providecommand keeps note-local definitions authoritative.
\usepackage{dsfont}
\providecommand{\ket}[1]{|#1\rangle}
\providecommand{\bra}[1]{\langle #1|}
\providecommand{\braket}[2]{\langle #1 | #2 \rangle}
\providecommand{\ketbra}[2]{|#1\rangle\!\langle #2|}
\providecommand{\Id}{\mathds{1}}
\providecommand{\one}{\mathds{1}}
\providecommand{\C}{\mathbb{C}}
\providecommand{\MN}[1]{M_{#1}(\C)}
\providecommand{\MD}{M_D(\C)}

% External referencing for paper sources.  The build helper writes sanitized
% auxiliary files under build/paper-aux/ and excludes duplicated source labels.
\usepackage{xr}

% Import a generated paper auxiliary file with a caller-supplied label prefix.
% The candidate paths support builds from docs/paper-gaps/, the repository root,
% and build/paper-aux/ itself.
\newcommand{\importpaperaux}[2]{%
  \IfFileExists{../../build/paper-aux/#2.aux}{%
    \externaldocument[#1]{../../build/paper-aux/#2}%
  }{%
    \IfFileExists{build/paper-aux/#2.aux}{%
      \externaldocument[#1]{build/paper-aux/#2}%
    }{%
      \IfFileExists{#2.aux}{%
        \externaldocument[#1]{#2}%
      }{}%
    }%
  }%
}

% General external reference macros
\newcommand{\paperref}[2]{\ref{#1-#2}}
\newcommand{\papereqref}[2]{(\ref{#1-#2})}

% CPSV16 source (1606.00608 / MPDO-22-12-17-2.tex) external document and shortcuts
\importpaperaux{cpsv16-}{1606.00608/MPDO-22-12-17-2-tnlean}
\newcommand{\cpsvref}[1]{\paperref{cpsv16}{#1}}
\newcommand{\cpsveqref}[1]{\papereqref{cpsv16}{#1}}

% Verdict marker: \gapnote{<kind>}{<status>}. Kind is the policy
% classification (clarification, local-correction, scope-restriction,
% unfaithful, false-source, open-gap); status is open, wip (elimination
% underway), resolved, or historical. Machine-read by the site index;
% prints a verdict line.
\newcommand{\gapnote}[2]{\par\noindent\textbf{Verdict:} #1 (#2).\par}


\title{Proof-Route Alignment for the Periodic Overlap Dichotomy
  (de las Cuevas--Cirac--Schuch--P\'erez-Garc\'ia)}
\author{}
\date{2026-06-04; revised 2026-08-02}

\begin{document}
\maketitle
\gapnote{scope-restriction}{open}

\section{Scope}

This note compares the formal proof of the periodic overlap dichotomy with the
Appendix~A argument of Ref.~\cite{DeLasCuevas2017Irreducible}. It identifies
the equivalent proof routes used in the formal development and the statements
that remain restricted relative to the source.\footnote{Formal files:
\path{TNLean/MPS/Periodic/Overlap/}. Local source:
\path{Papers/1708.00029/main.tex}. Paper-level umbrella:
\href{https://github.com/LionSR/TNLean/issues/619}{issue \#619};
proposition-level subtracking:
\href{https://github.com/LionSR/TNLean/issues/81}{issue \#81}. The Case-3
contraction and phase construction was tracked by
\href{https://github.com/LionSR/TNLean/issues/873}{issue \#873} and is now
formalized by \leanid{MPSTensor.sectorTensor_proportional_of_blockedMatch}.}

The source proposition is ``equal-or-orthogonal-generalized''.\footnote{The
statement appears at lines 589--609 and its proof appears at lines 903--1118
of Ref.~\cite{DeLasCuevas2017Irreducible}.} The differences recorded below do
not introduce hidden hypotheses. Each formal statement is faithful to the
source or explicitly restricts its scope, and each substituted proof route is
mathematically equivalent to the source route.

\section{Positive chain lengths}

The formal periodic-overlap argument states its matrix-product-vector
equalities for $N\geq 1$. In particular,
\leanid{MPSTensor.peripheralProportionalCase_periodicFT_of_sameMPV₂Pos} assumes equality
at every positive length. This assumption suffices because the overlap
argument is asymptotic and uses only a tail of positive lengths. The empty word
contains only a bond-dimension trace and does not enter theorem
\emph{thm:bd}, lines 613--623, of
Ref.~\cite{DeLasCuevas2017Irreducible}. The formal development also propagates
nonzero cyclic-sector dimensions directly from the corner-algebra
equivalences and the cyclic projector shift, without using an empty-word trace
identity.

\section{Different-period decay through the peripheral spectrum}

\paragraph{Source route.}

For periods $m_a\neq m_b$, Appendix~A, lines 917--950, of
Ref.~\cite{DeLasCuevas2017Irreducible} blocks at
$p=\operatorname{lcm}(m_a,m_b)$ and applies Lemma~\emph{bdcf}. The tensors
$P_uA^{(p)}$ and $Q_vB^{(p)}$ are non-repeated normal tensors. Applying the
one-site translation operator to $Q_vB^{(p)}$ and $Q_{v+1}B^{(p)}$ gives a
contradiction unless the overlap decays:
\begin{align}
    \lim_{N\to\infty}
    \left|\braket{V_N(A)}{V_N(B)}\right|
    &=0.
    \label{eq:dccsp17_periodic_overlap_route_alignment_different_period_decay}
\end{align}

\paragraph{\Lean{} route.}

For a tensor $A=(A^i)_i$, define its transfer map and adjoint by
\begin{align}
    \mathcal E_A(Y)
    &=\sum_i A^iY(A^i)^\dagger,
    \label{eq:dccsp17_periodic_overlap_route_alignment_transfer_map}\\
    \mathcal E_A^*(Y)
    &=\sum_i (A^i)^\dagger YA^i.
    \label{eq:dccsp17_periodic_overlap_route_alignment_transfer_adjoint}
\end{align}
Suppose $D_a=D_b$ and
\begin{align}
    B^i
    &=\zeta XA^iX^{-1}.
    \label{eq:dccsp17_periodic_overlap_route_alignment_gauge_phase_relation}
\end{align}
Then the transfer maps satisfy
\begin{align}
    \mathcal E_B(Y)
    &=\zeta\bar\zeta
      X\mathcal E_A\!\left(X^{-1}YX^{-\dagger}\right)X^\dagger.
    \label{eq:dccsp17_periodic_overlap_route_alignment_transfer_similarity}
\end{align}
The uniqueness of the peripheral eigenvalue for an irreducible completely
positive map, Theorem~6.3 of Ref.~\cite{Wolf2012Quantum}, forces
$|\zeta|=1$. Hence $\mathcal E_A$ and $\mathcal E_B$ have the same peripheral
spectrum. An $m$-periodic tensor has peripheral spectrum
\begin{align}
    \operatorname{spec}_{\mathrm{per}}(\mathcal E_A)
    &=\left\{e^{2\pi ir/m}:0\leq r<m\right\}.
    \label{eq:dccsp17_periodic_overlap_route_alignment_periodic_spectrum}
\end{align}
Equality of the peripheral spectra therefore implies $m_a=m_b$. Different
periods exclude gauge-phase equivalence, so the irreducible trace-preserving
equal-or-orthogonal dichotomy gives
(\ref{eq:dccsp17_periodic_overlap_route_alignment_different_period_decay}).
This spectral argument replaces the least-common-multiple blocking and
translation argument of the source and is complete.\footnote{Formal
declaration:
\leanid{MPSTensor.periodicOverlap_tendsto_zero_of_ne_period} in
\path{DifferentPeriod.lean}.}

\section{Sector non-repetition from the blocked fixed-point structure}

\paragraph{Source argument.}

Lemma~\emph{bdcf}, lines 404--423, of
Ref.~\cite{DeLasCuevas2017Irreducible} considers an irreducible periodic block
whose transfer map has peripheral spectrum
\begin{align}
    \operatorname{spec}_{\mathrm{per}}(\mathcal E_A)
    &=\{\omega^r:0\leq r<m\},
    \label{eq:dccsp17_periodic_overlap_route_alignment_block_peripheral_spectrum}\\
    \omega
    &=e^{2\pi i/m}.
    \label{eq:dccsp17_periodic_overlap_route_alignment_primitive_root}
\end{align}
The blocked map $\mathcal E_A^m$ has $1$ as its only modulus-one eigenvalue,
with multiplicity $m$, and its fixed-point space is
\begin{align}
    \operatorname{Fix}(\mathcal E_A^m)
    &=\operatorname{span}
      \{P_u\Lambda_AP_u:0\leq u<m\}.
    \label{eq:dccsp17_periodic_overlap_route_alignment_block_fixed_points}
\end{align}
The adjoint fixed-point space is
\begin{align}
    \operatorname{Fix}(\mathcal E_A^{*m})
    &=\operatorname{span}
      \{P_u:0\leq u<m\}.
    \label{eq:dccsp17_periodic_overlap_route_alignment_adjoint_fixed_points}
\end{align}

Assume, for distinct sectors $u\neq v$, that there is a gauge-phase
equivalence
\begin{align}
    C_u^{\mathbf i}
    &=e^{i\xi}UC_v^{\mathbf i}U^\dagger,
    \label{eq:dccsp17_periodic_overlap_route_alignment_sector_gauge_phase}\\
    U
    &=P_uUP_v,
    \label{eq:dccsp17_periodic_overlap_route_alignment_off_diagonal_support}\\
    UU^\dagger
    &=P_u,
    \label{eq:dccsp17_periodic_overlap_route_alignment_left_support_isometry}\\
    U^\dagger U
    &=P_v.
    \label{eq:dccsp17_periodic_overlap_route_alignment_right_support_isometry}
\end{align}
Using the normalization
$\sum_{\mathbf i}C_v^{\mathbf i}C_v^{\mathbf i\dagger}=P_v$, one obtains
\begin{align}
    \mathcal E_A^m(U)
    &=\sum_{\mathbf i}
      C_u^{\mathbf i}U C_v^{\mathbf i\dagger},
    \label{eq:dccsp17_periodic_overlap_route_alignment_off_diagonal_transfer}\\
    &=e^{i\xi}U
      \sum_{\mathbf i}
      C_v^{\mathbf i}C_v^{\mathbf i\dagger},
    \label{eq:dccsp17_periodic_overlap_route_alignment_insert_sector_gauge}\\
    &=e^{i\xi}U.
    \label{eq:dccsp17_periodic_overlap_route_alignment_off_diagonal_eigenvector}
\end{align}
Thus $U$ is a modulus-one eigenvector of $\mathcal E_A^m$. However,
(\ref{eq:dccsp17_periodic_overlap_route_alignment_block_fixed_points}) contains
only diagonal sector elements, whereas
(\ref{eq:dccsp17_periodic_overlap_route_alignment_off_diagonal_support}) is
off-diagonal. Since $U\neq 0$, this is a contradiction.

\paragraph{\Lean{} status.}

The formal statement proves the same non-repetition result: distinct cyclic
sectors with $P_uP_v=0$ are not gauge-phase equivalent.\footnote{Formal
declaration:
\path{MPSTensor.not_gaugePhaseEquiv_of_orthogonal_cyclicSector_traces} in
\path{SelfOverlap.lean}.} The proof follows the spectral route above. From a
gauge-phase equivalence between two compressed sectors, it constructs a
nonzero off-diagonal eigenvector of $\mathcal E_A^m$ using the sector
corner-letter identity and rectangular support isometries for the compressed
corners. The existing off-diagonal vanishing lemma then gives the
contradiction. The scalar normalization uses the same Perron--Frobenius
eigenvalue-uniqueness input as the different-period proof.\footnote{Formal
declarations:
\leanid{MPSTensor.exists_offDiag_eigenvector_of_gaugePhase_cast_left},
\leanid{MPSTensor.offDiag_eigenvector_eq_zero_of_isPeriodic}, and
\path{MPSTensor.period_eq_of_gaugePhaseEquiv_of_isPeriodic}.}

\paragraph{Signature correction.}

Periodicity is essential. The conclusion is false under only orthogonal
projectors $P_k$ summing to the identity, the adjoint shift
$\mathcal E_A^*(P_{k+1})=P_k$, blocked-letter commutation, the trace formula,
and orthogonality. A reducible tensor $A=B\oplus B$, where $B$ is an
irreducible period-$m$ block, admits these cyclic-projection data but has
gauge-phase-equivalent distinct sectors. The spectral proof requires $A$ to
be a periodic irreducible block, because periodicity supplies irreducibility,
the prescribed peripheral spectrum, and primitive sectors.

The hypothesis \path{hP : IsPeriodic m A} was therefore added to the lemma
signature. It is supplied at the unique call site
\path{MPSTensor.sectorBlocks_not_gaugePhaseEquiv_of_ne}. With this hypothesis and the
corner-isometry data, the spectral proof is formalized.

\section{Sector matching: propagation and phase construction}

\paragraph{Source route.}

Appendix~A, lines 961--1117, of
Ref.~\cite{DeLasCuevas2017Irreducible} starts from the blocked identity
\emph{eq:Nm}, lines 978--984. Applying the translation operator $T^l$ and
Theorem~2.10 of Ref.~\cite{Cirac2016MPDO} produces, at each
 offset, a phase
and a unitary
\begin{align}
    U_{\widetilde v+l}
    &=P_{\widetilde u+l}U_{\widetilde v+l}Q_{\widetilde v+l}.
    \label{eq:dccsp17_periodic_overlap_route_alignment_offset_unitary}
\end{align}
The offset is therefore constant:
\begin{align}
    q
    &=v-u.
    \label{eq:dccsp17_periodic_overlap_route_alignment_constant_offset}
\end{align}

Each blocked product $F_u$ becomes injective after $N_0$ repetitions and has a
right inverse $\Omega_u$. Contracting the concatenated products with these
right inverses recovers the single-site relation
\begin{align}
    A_u^i
    &=\kappa_v e^{i\eta/m}B_v^i.
    \label{eq:dccsp17_periodic_overlap_route_alignment_site_proportionality}
\end{align}
The phase calculation uses the source identity \emph{eq:prodkappaprop}:
\begin{align}
    \prod_v\kappa_v
    &=1,
    \label{eq:dccsp17_periodic_overlap_route_alignment_kappa_product}\\
    |\kappa_v|
    &=1,
    \label{eq:dccsp17_periodic_overlap_route_alignment_kappa_unit_modulus}\\
    \kappa_v
    &=e^{i\theta_v},
    \label{eq:dccsp17_periodic_overlap_route_alignment_kappa_phase}\\
    \sum_v\theta_v
    &=0.
    \label{eq:dccsp17_periodic_overlap_route_alignment_theta_sum}
\end{align}
Choose phases $\phi_v$ satisfying
\begin{align}
    \theta_v
    &=\phi_v-\phi_{v+1}.
    \label{eq:dccsp17_periodic_overlap_route_alignment_phase_coboundary}
\end{align}
The sector phases then telescope into the global phase and unitary
\begin{align}
    \xi
    &=\eta/m,
    \label{eq:dccsp17_periodic_overlap_route_alignment_global_phase}\\
    U
    &=\sum_u
      e^{i\phi_{u+q}}P_uU_{u+q}Q_{u+q},
    \label{eq:dccsp17_periodic_overlap_route_alignment_global_unitary}\\
    A^i
    &=e^{i\xi}UB^iU^\dagger.
    \label{eq:dccsp17_periodic_overlap_route_alignment_global_gauge}
\end{align}

\paragraph{\Lean{} status.}

The formal development follows the same two stages: the $T^l$ and
Theorem~2.10 propagation staircase, followed by the $\Omega_u$ contraction and
the $\kappa/\theta/\phi$ phase construction.\footnote{Formal declarations:
\leanid{MPSTensor.sectorMatch_propagation} in
\path{SectorMatch/Propagation.lean},
\leanid{MPSTensor.sectorTensor_proportional_of_blockedMatch} in
\path{SectorMatch/Contraction.lean}, and
\leanid{MPSTensor.periodicOverlap_gaugeEquiv_of_sector_match} in
\path{SectorMatch/Consequences.lean}.} The top-level theorem invokes these
component lemmas. The contraction leaf
\leanid{MPSTensor.sectorTensor_proportional_of_blockedMatch}, the global unitary gauge,
and the repeated-blocks theorem are proved.

The docstrings now cite the source equation identifiers and line ranges rather
than published-version identifiers such as ``A.8'' and ``A.14--A.18'', which
do not occur in the local source. They also state explicitly that the
$\kappa/\theta/\phi$ construction is essential.

\subsection{Single-site inputs}

The formal single-site off-diagonal decomposition follows from the cyclic
adjoint-shift relation.\footnote{Formal declarations:
\leanid{MPSTensor.offDiag_shift_of_adjoint_cyclic_shift},
\leanid{MPSTensor.eq_sum_offDiag_of_adjoint_cyclic_shift},
\leanid{MPSTensor.IsCyclicSectorDecomp.offDiag_shift}, and
\leanid{MPSTensor.IsCyclicSectorDecomp.eq_sum_offDiag}.} With the repository's inverse
cyclic indexing convention, it gives
\begin{align}
    P_{u+1}A^i
    &=A^iP_u,
    \label{eq:dccsp17_periodic_overlap_route_alignment_single_site_shift}\\
    A^i
    &=\sum_u P_{u+1}A^iP_u.
    \label{eq:dccsp17_periodic_overlap_route_alignment_single_site_decomposition}
\end{align}
These identities are the inverse-indexed form of source equation
\emph{eq:Aoffdiag}, lines 335--339, of
Ref.~\cite{DeLasCuevas2017Irreducible}. The later source equation
\emph{eq:Auprop}, line 1014, defines the sector letters by
\begin{align}
    A_u^i
    &=P_uA^iP_{u+1}.
    \label{eq:dccsp17_periodic_overlap_route_alignment_sector_letter}
\end{align}

The cyclic-sector construction also records a letter-level corner identity
and a rectangular support isometry for each corner.\footnote{Formal
declarations:
\leanid{MPSTensor.exists_cyclic_sector_decomp_with_letter_and_isometry_after_blocking_of_isPeriodic}
and
\leanid{MPSTensor.exists_cyclic_sector_decomp_with_letter_after_blocking_of_isPeriodic}.}
If $\varphi_u$ is the compression isomorphism from the compressed sector
algebra onto $P_uM_DP_u$, then
\begin{align}
    \varphi_u(C_u^{\mathbf i})
    &=P_u(A^{(m)})^{\mathbf i}P_u.
    \label{eq:dccsp17_periodic_overlap_route_alignment_corner_letter}
\end{align}
This identity formalizes the blocked corner-letter relation in Lemma
\emph{bdcf} and source equation \emph{Cu} of
Ref.~\cite{DeLasCuevas2017Irreducible}. Together with the support isometry, it
moves the compressed mixed-transfer eigenvector into the ambient corner and
supplies the off-diagonal modulus-one eigenvector used in the spectral
contradiction.

Define the sector overlap by
\begin{align}
    O_{C,D}(N)
    &:=\braket{V_N(C)}{V_N(D)}.
    \label{eq:dccsp17_periodic_overlap_route_alignment_sector_overlap}
\end{align}
The one-site sector-overlap translation is also formalized.\footnote{Formal
declaration:
\leanid{MPSTensor.sectorOverlap_succ_eq_of_cyclicSectorDecomp}.} For every $N>0$, it
proves
\begin{align}
    O_{A_{u+1},B_{v+1}}(N)
    &=O_{A_u,B_v}(N).
    \label{eq:dccsp17_periodic_overlap_route_alignment_overlap_translation}
\end{align}
This is the overlap-sum form of the translation step at lines 985--1002 of
Ref.~\cite{DeLasCuevas2017Irreducible}. It applies the single-site relations
inside the paired trace sum using the inverse cyclic indexing convention.

\subsection{Formalized Case-3 contraction}

The formal contraction and gauge construction has two parts.
\begin{itemize}
    \item The $m$-factor cyclic $\Omega_u$ contraction starts from common right
    inverses for the injective tensors $F_u$ and contracts the cyclic product
    at lines 1041--1068 of Ref.~\cite{DeLasCuevas2017Irreducible}. This yields
    the uniform product-tensor identity \emph{eq:resultprop}, with the $\eta$
    phases retained. The displayed inverse list at lines 1049--1052 contains
    only $m-1$ factors, but the displayed concatenation, the $(mN_0+1)$
    repetition count, and the definition of $\eta_v$ require one inverse for
    each of the $m$ inserted $F$-factors.

    \item The normalization and gauge step extracts scalars $\kappa_v$,
    proves (\ref{eq:dccsp17_periodic_overlap_route_alignment_kappa_product})
    and
    (\ref{eq:dccsp17_periodic_overlap_route_alignment_kappa_unit_modulus}),
    applies the finite-cycle coboundary lemma
    (\ref{eq:dccsp17_periodic_overlap_route_alignment_phase_coboundary}), and
    assembles the unitary
    (\ref{eq:dccsp17_periodic_overlap_route_alignment_global_unitary}).
\end{itemize}
These steps discharge the contraction obligation for the repeated-blocks
theorem.\footnote{Formal declaration:
\leanid{MPSTensor.sectorTensor_proportional_of_blockedMatch} in
\path{SectorMatch/Contraction.lean}.} Tensor-product scalar extraction and
product-one normalization are isolated in
\leanid{TNLean.Algebra.PiTensorProductPhase.exists_kappa_of_piTensorProduct_eq_smul} and
\leanid{TNLean.Algebra.PiTensorProductPhase.exists_kappa_product_one_of_piTensorProduct_eq_root_smul}.
The finite-cycle phase choice is isolated in
\leanid{TNLean.Algebra.exists_fin_complex_unit_cyclic_coboundary_shift_of_prod_eq_one}.
The one-step transport theorem follows from the single-site overlap theorem,
the overlap-to-gauge theorem, and sector primitivity.\footnote{Formal
declarations:
\path{MPSTensor.sectorGaugePhaseEquiv_succ_of_cyclicTransport} and
\leanid{MPSTensor.gaugePhaseEquiv_of_overlap_norm_tendsto_one_of_irreducible_TP}.}

\subsection{Corner-letter building blocks}

The formal corner-transition tensors and their repeated products are
\begin{align}
    A_u^i
    &=P_uA^iP_{u+1},
    \label{eq:dccsp17_periodic_overlap_route_alignment_corner_transition}\\
    F_u^{\mathbf i}
    &=A_u^{i_1}A_{u+1}^{i_2}\cdots.
    \label{eq:dccsp17_periodic_overlap_route_alignment_corner_product}
\end{align}
They correspond to source equations \emph{eq:Auprop} and \emph{eq:Fu} of
Ref.~\cite{DeLasCuevas2017Irreducible} and are defined as
\leanid{MPSTensor.cornerLetter} and \leanid{MPSTensor.cornerProd}. Their cyclic
concatenation law is
\begin{align}
    F_u^{\mathbf i_1\mathbf i_2}
    &=F_u^{\mathbf i_1}
      F_{u+\ell_1}^{\mathbf i_2}.
    \label{eq:dccsp17_periodic_overlap_route_alignment_corner_concatenation}
\end{align}
The declaration is \leanid{MPSTensor.cornerProd_append}. The junction
projector is absorbed by idempotency, so this law requires neither injectivity
nor normality. These are the $A_u^i$ and $F_u$ objects used in the cyclic
contraction. The formal right-inverse contraction to source equation
\emph{eq:resultprop} and the global unitary assembly are completed in
\leanid{MPSTensor.sectorTensor_proportional_of_blockedMatch}.

\section{Scope restriction: independence without spanning}

The source proposition and Lemma~\emph{Lem1t}, lines 511--519, imply that for
$N\geq N_0$ the nonzero vectors in
\begin{align}
    \left\{\ket{V_N(A_j)}:j\in J\right\}
    \label{eq:dccsp17_periodic_overlap_route_alignment_periodic_vector_family}
\end{align}
are linearly independent and span $\ket{V_N(A)}$; see lines 604--608 of
Ref.~\cite{DeLasCuevas2017Irreducible}. The spanning assertion justifies the
source phrase ``basis of periodic vectors'' at line 611.

The formal theorem states only eventual linear independence.\footnote{Formal
declaration:
\leanid{MPSTensor.periodicBasis_eventuallyLinearlyIndependent} in
\path{Dichotomy.lean}. The corresponding blueprint entry is in
\path{blueprint/src/chapter/ch22_periodic_ft_overlap_sector_match_and_consequences.tex}.}
It encodes the basis assumptions directly as the pairwise non-repetition
hypothesis \path{hNonrep} and the common-period divisibility hypothesis
\path{hDiv}. It neither derives them from a basis of periodic tensors nor
uses the almost-orthonormal-implies-independent Lemma~\emph{Lem1t}. Restoring
the spanning clause and routing the proof through \emph{Lem1t} was classified
under \href{https://github.com/LionSR/TNLean/issues/81}{issue \#81}. The
independence clause is formalized, but no unrestricted spanning theorem is
currently attributed to the source consequence.

\section{Multiplicity-bearing overlap theorem and Perron transport}

\subsection{Source irreducible forms}

The proportional Fundamental Theorem uses the irreducible forms
\emph{bdnr} and \emph{Bbdnr}, lines 286--305 and 575--585, of
Ref.~\cite{DeLasCuevas2017Irreducible}:
\begin{align}
    A^i
    &=\bigoplus_j(R_j\otimes A_j^i),
    \label{eq:dccsp17_periodic_overlap_route_alignment_irreducible_form_A}\\
    B^i
    &=\bigoplus_k(S_k\otimes B_k^i).
    \label{eq:dccsp17_periodic_overlap_route_alignment_irreducible_form_B}
\end{align}
The diagonal matrices $R_j$ and $S_k$ retain the multiplicities and weights of
repeated copies of each periodic basis tensor.

\subsection{Multiplicity-bearing normalized theorem}

The declaration
\path{PeriodicOverlapHypothesis.ofSectorDecompositions} uses literal sector
decompositions
\begin{align}
    P^i
    &=\bigoplus_{j,q}\mu_{j,q}A_j^i,
    \label{eq:dccsp17_periodic_overlap_route_alignment_literal_decomposition_P}\\
    Q^i
    &=\bigoplus_{k,r}\nu_{k,r}B_k^i.
    \label{eq:dccsp17_periodic_overlap_route_alignment_literal_decomposition_Q}
\end{align}
Every complex weight in these decompositions is nonzero. This flattened
notation is the coordinate form of
(\ref{eq:dccsp17_periodic_overlap_route_alignment_irreducible_form_A}) and
(\ref{eq:dccsp17_periodic_overlap_route_alignment_irreducible_form_B}), with
\begin{align}
    \tr(R_j^N)
    &=\sum_q\mu_{j,q}^N,
    \label{eq:dccsp17_periodic_overlap_route_alignment_R_power_sum}\\
    \tr(S_k^N)
    &=\sum_r\nu_{k,r}^N.
    \label{eq:dccsp17_periodic_overlap_route_alignment_S_power_sum}
\end{align}
The theorem assumes proportionality only at positive lengths, permits
different total bond dimensions, and allows the proportionality scalar to
depend on the length. It therefore removes the multiplicity,
literal-representation, and empty-word restrictions of the earlier scalar
theorem.

The proof expands the contradiction stated tersely at line 631 of
Ref.~\cite{DeLasCuevas2017Irreducible}. It partitions one periodic basis by
mixed-overlap decay, restricts to a common-period subsequence, and compares
grouped coefficients in an eventually linearly independent family indexed by
a coproduct. Linear independence yields
\begin{align}
    \tr(R_{j_0}^{pN})
    &=0
    \qquad\text{for every sufficiently large }N.
    \label{eq:dccsp17_periodic_overlap_route_alignment_vanishing_power_sum}
\end{align}
Geometric extrapolation applied to the nonzero numbers $\mu_{j_0,q}^p$ then
forces the exponent-zero power sum to vanish:
\begin{align}
    \sum_q(\mu_{j_0,q}^p)^0
    &=0.
    \label{eq:dccsp17_periodic_overlap_route_alignment_zero_exponent_contradiction}
\end{align}
The left-hand side is the positive multiplicity of the sector, which is a
contradiction. The proof never treats repeated copies of $A_{j_0}$ as
independent vectors.

The earlier declaration
\path{PeriodicOverlapHypothesis.ofIsIrreducibleForm} remains the specialization
with one weight per listed block. Its hypotheses still assert
matrix-product-vector equality at the empty word and may specify an ambient
tensor only through that equality. These are restrictions of the older
formulation, not of the literal multiplicity-bearing theorem.

\subsection{Pure normalization theorem}

The predicate \path{IsSpectrallyPeriodic} records the source assumptions on a
basis block: irreducibility, spectral radius one, positive period, and the
prescribed roots-of-unity peripheral spectrum. The declaration
\path{IsSpectrallyPeriodic.exists_isPeriodic_tpGauge} applies the positive
adjoint Perron eigenmatrix to obtain a left-canonical periodic block by pure
similarity, following lines 313--332 of
Ref.~\cite{DeLasCuevas2017Irreducible}. The Perron eigenvalue equals one, so
the construction does not rescale the block. Similarity preserves the
peripheral spectrum. Consequently, the normalization leaves every
matrix-product vector unchanged and introduces no additional power into the
multiplicity coefficients.

\subsection{Sectorwise normalization and return to the original blocks}

The declaration \path{SectorDecomposition.replaceBasis} replaces only the
representative tensors. It preserves the sector index set, bond dimensions,
multiplicities, and weights. If $X_j$ is the pure Perron gauge for the $j$th
representative, then
\path{SectorDecomposition.gaugeEquiv_toTensor_replaceBasis} assembles
\begin{align}
    X
    &=\bigoplus_j(I_{r_j}\otimes X_j).
    \label{eq:dccsp17_periodic_overlap_route_alignment_sectorwise_gauge}
\end{align}
Thus each block transforms as
\begin{align}
    \mu_{j,q}A_j^i
    &\longmapsto
      \mu_{j,q}X_jA_j^iX_j^{-1},
    \label{eq:dccsp17_periodic_overlap_route_alignment_weight_preservation}
\end{align}
without changing $\mu_{j,q}$. The declaration
\path{SectorDecomposition.exists_isPeriodic_replaceBasis} applies this
normalization simultaneously to all spectrally periodic representatives.

Pure similarities preserve every closed-chain coefficient and hence every
mixed overlap. The declaration
\path{PeriodicOverlapHypothesis.of_gaugeEquiv} uses this equality to preserve
overlap decay and non-decay, and it composes the pure gauges with any
repeated-block relation obtained after normalization. The same transitivity
argument preserves pairwise non-repetition.

Finally,
\path{PeriodicOverlapHypothesis.ofSpectrallyPeriodicSectorDecompositions}
normalizes both literal multiplicity-bearing decompositions, transports their
positive-length proportionality to the normalized tensors, applies the
normalized multiplicity theorem, and returns the complete periodic-overlap
hypothesis to the original spectral-radius-one blocks. This completes the
overlap-hypothesis step tracked by
\href{https://github.com/LionSR/TNLean/issues/5342}{issue \#5342}.

The declarations \leanid{MPSTensor.fundamentalTheorem_periodic_proportional} and
\leanid{MPSTensor.fundamentalTheorem_periodic_proportional_of_isPeriodic} start after
this step. They assume the periodic-overlap hypothesis, or its non-decaying
partner fields, and prove finite block matching. Together with the spectrally
periodic multiplicity theorem, they yield block matching for literal
irreducible forms.

\subsection{Equality of the matched periods}

The matching clause of Theorem~3.4, lines 616--620, of
Ref.~\cite{DeLasCuevas2017Irreducible} asserts that the unique matched block
carries the same period. The source justifies this by the last clause of the
generalized equal-or-orthogonal proposition, lines 602--604: the repeated-block
relation ``can only happen if $m_a=m_b$''.

The formal argument is the spectral one already used for different-period
decay. If $A^i=\xi YB^iY^{-1}$ with $|\xi|=1$, then
\begin{align}
    \mathcal E_A(Z)
    &=Y\mathcal E_B\!\left(Y^{-1}ZY^{-\dagger}\right)Y^\dagger,
    \label{eq:dccsp17_periodic_overlap_route_alignment_repeated_transfer_congruence}
\end{align}
because the unit-modulus scalar contributes the factor $\xi\bar\xi=1$.
Conjugation by an invertible linear map preserves the spectrum, so the
peripheral spectra of $\mathcal E_A$ and $\mathcal E_B$ agree.\footnote{Formal
declaration:
\leanid{MPSTensor.RepeatedBlocks.peripheralEigenvalues_transferMap_eq}.} For
spectrally periodic blocks those spectra are the sets of roots of unity of
orders $m$ and $n$, and
(\ref{eq:dccsp17_periodic_overlap_route_alignment_periodic_spectrum}) then
forces $m=n$ by mutual divisibility of the two orders.\footnote{Formal
declarations: \leanid{MPSTensor.period_eq_of_setOf_pow_eq_one} and
\leanid{MPSTensor.IsSpectrallyPeriodic.period_eq_of_hetRepeatedBlocks}.}

The source-labelled theorem is therefore formalized. It combines the
spectrally periodic overlap theorem with the finite matching argument and
records the matched periods in the conclusion.\footnote{Formal declaration:
\leanid{MPSTensor.fundamentalTheorem_periodic_proportional_sectorDecomposition},
over the witness \path{MPSTensor.PeriodicBasisMatchingWitness}. Blueprint
entry: \emph{thm:periodic\_ft\_proportional}.} Its hypotheses are the source
hypotheses: multiplicity-bearing irreducible forms with nonzero diagonal
entries, spectrally periodic and pairwise non-repeated bases of periodic
tensors, and nonzero proportionality of the generated matrix-product vectors
at every positive length.

\subsection{Per-length proportionality and the blockwise phase}

The hypothesis of Theorem~3.4, lines 613--623, of
Ref.~\cite{DeLasCuevas2017Irreducible} is projective proportionality separately
at each positive length:
\begin{align}
    \forall N\geq1,\ \exists c_N\neq0,\qquad
    \ket{V_N(A)}
    &=c_N\ket{V_N(B)}.
    \label{eq:dccsp17_periodic_overlap_route_alignment_per_length_proportionality}
\end{align}
No scalar independent of $N$ is assumed, and none may be introduced. The
one-dimensional blocks $B^0=1$ and $A^0=e^{i\theta}$ satisfy
\begin{align}
    \ket{V_N(A)}
    &=e^{iN\theta}\ket{V_N(B)},
    \label{eq:dccsp17_periodic_overlap_route_alignment_phase_power_example}
\end{align}
so a length-independent scalar is false unless the phase is trivial. The
formal statements accordingly quantify the scalar over the length, and the
repeated-block consequence records the scalar in the exponential form
\begin{align}
    \ket{V_N(A_j)}
    &=\zeta_j^{\,N}\ket{V_N(B_{\pi(j)})},
    \qquad|\zeta_j|=1,
    \label{eq:dccsp17_periodic_overlap_route_alignment_blockwise_phase_power}
\end{align}
which follows from the repeated-block relation \emph{eq:rep}, lines 276--284,
by cancelling the similarity inside the closed-chain trace.\footnote{Formal
declaration:
\leanid{MPSTensor.HetRepeatedBlocks.exists_unit_phase_power_mpv}.}

The conclusion of Theorem~3.4 is blockwise. Each matched pair carries its own
phase $\xi_j$ in
\begin{align}
    A_j^i
    &=e^{i\xi_j}Y_jB_{\pi(j)}^iY_j^{-1},
    \label{eq:dccsp17_periodic_overlap_route_alignment_blockwise_relation}
\end{align}
and distinct matched pairs may carry distinct phases. The source states at
lines 625--627 that the theorem relates only the bases of periodic tensors and
not the assembled tensors themselves, so no single phase for the assembled
tensors is available at this stage. The formal phase-absorption statement is
therefore also blockwise: for one matched pair it produces a unit-modulus
$\xi$ with $\ket{V_N(\xi A_j)}=\ket{V_N(B_{\pi(j)})}$ for every $N\geq1$,
matching the phase absorption used at lines 302--305 and again at lines
667--671.\footnote{Formal declaration:
\leanid{MPSTensor.HetRepeatedBlocks.exists_phase_rescaling_sameMPV₂Pos}.} Its
scope restriction is recorded on the declaration itself. Relating the
assembled tensors requires the multiplicities and is possible only in the
equal case, theorem \emph{thm:bdequal}, lines 643--658.

\subsection{Remaining elimination plan}

Extending the irreducible-form theorem to arbitrary input tensors through
Proposition~2.1 of Ref.~\cite{DeLasCuevas2017Irreducible} requires the
positive-length, possibly lower-dimensional comparison at lines 266--271 of
the same source. That extension is not part of Theorem~3.4. The completed
overlap-hypothesis step is tracked by
\href{https://github.com/LionSR/TNLean/issues/5342}{issue \#5342}; the
subsequent multiplicity-power extraction is tracked by
\href{https://github.com/LionSR/TNLean/issues/5344}{issue \#5344}.

\section{The equal case with trace-preserving blocks: restriction eliminated}

Theorem~3.8 of Ref.~\cite{DeLasCuevas2017Irreducible}, lines 643--693, assumes
that the two tensors are in irreducible form, that is, that every block map is
irreducible of spectral radius one. The equal-case argument runs in the
normalized orientation, since the periodic overlap dichotomy and the vanishing
of the matrix-product vector of a periodic block at lengths that its period
does not divide are both available only there. It therefore first produces the
statement for blocks whose maps are trace preserving, the normalization of
Eq.~\emph{unital}, line 316.\footnote{Formal declaration:
\leanid{MPSTensor.fundamentalTheorem_periodic_equalCase_sectorDecomposition}.}

That hypothesis is worth placing precisely, since the source's irreducible form
II is a different and stronger condition: form II, at lines 320--332, asks for
trace-preserving block maps \emph{and} for the fixed point of each block map to
be diagonal, the latter obtained by the further conjugation
$A''_j = U_j^\dagger A'_j U_j$. Only the first of the two is assumed, so the
normalized statement is stronger than the form-II statement.

The passage that the source asserts at lines 330--332 is carried out, so the
unrestricted source statement is
formalized.\footnote{Formal declaration:
\leanid{MPSTensor.fundamentalTheorem_periodic_equalCase_irreducibleForm}.} Only
the unitalization of Eq.~\emph{unital} is needed: each representative is
conjugated by the square root of the positive-definite fixed point of the
adjoint block map. Being blockwise, it leaves every multiplicity entry and
every multiplicity power sum untouched, and the multiplicity gauge, which is
scalar on the bond space of each multiplicity copy, commutes with it and
returns unchanged. Only the similarity between the two tensors is conjugated.
The second half of the form-II construction, the diagonalization of the fixed
point, is not needed.

The conclusion asserts that the similarity is invertible. The source remark at
lines 330--332 and 625 that in irreducible form II the similarity is unitary is
a separate statement and is not asserted here.

\section{Scope restriction: the periods of an irreducible-form block are
    supplied}

The source defines irreducible form at line 261 by two conditions on each
block: its transfer map is an irreducible completely positive map, and its
spectral radius is one. The period $m_j$ is then \emph{derived}: at lines
257--258, citing Ref.~\cite{Wolf2012Quantum}, the source records that the
peripheral spectrum of such a map is exactly the set of $m_j$-th roots of unity
for some $m_j$, and calls the block periodic of period $m_j$ on that basis.
Theorem~3.8 accordingly mentions $m_j$ in its conclusion, not among its
hypotheses.

The formal equal-case theorem\footnote{Formal declaration:
\leanid{MPSTensor.fundamentalTheorem_periodic_equalCase_irreducibleForm}.}
instead receives the periods as input, together with witnesses that each block
is spectrally periodic of the given period.\footnote{Formal declaration:
\leanid{MPSTensor.IsSpectrallyPeriodic}.} That predicate carries four fields,
of which the source's definition supplies only the first two; the remaining
two, positivity of $m$ and the identification of the peripheral spectrum with
the $m$-th roots of unity, are the derived content. A reader holding exactly
the source's hypotheses therefore cannot instantiate the formal theorem, and
the same applies to the proportional case, Theorem~3.4.\footnote{Formal
declaration:
\leanid{MPSTensor.fundamentalTheorem_periodic_proportional_sectorDecomposition}.}

\subsection{Elimination plan}

What is missing is one existence statement: for a family $A$ with irreducible
transfer map of spectral radius one, there is an $m$ with
$\operatorname{Spec}_{\mathrm{per}}(\E_A) = \{\mu : \mu^m = 1\}$ and $m>0$. The
ingredients are present.

\begin{enumerate}
    \item The unitalization of Eq.~\emph{unital}, line 316, is already
          formalized from exactly the source's two conditions, and it preserves
          irreducibility and the peripheral spectrum.\footnote{Formal
          declaration:
          \leanid{MPSTensor.exists_tpGauge_of_irreducible_spectralRadius_one}.}
    \item The cyclic structure of the peripheral spectrum of a unital
          irreducible finite-Kraus map with a positive-definite adjoint fixed
          point is available in the companion library, in the form
          $\operatorname{Spec}_{\mathrm{per}} = \{\gamma^k : k<m\}$ with
          $\gamma$ a primitive $m$-th root of unity.\footnote{Formal
          declaration:
          \leanid{PeripheralSpectrum.peripheral_eigenvalues_cyclic_structure}.}
          It carries its own scope restriction, recorded in that library's
          note on Wolf Theorem 6.6.
    \item For a primitive $m$-th root $\gamma$ over $\C$ the two descriptions
          agree, $\{\gamma^k : k<m\} = \{\mu : \mu^m=1\}$.
\end{enumerate}

The assembly has to pass between the unital and trace-preserving orientations,
which exchange a Kraus family with its adjoint family and conjugate the
peripheral spectrum; the same conjugation is already carried out for the
transported statement in \path{SelfOverlapSetup.lean}. Once the existence
statement is available, both the proportional and the equal case can be
restated over the source's two conditions with the periods produced rather than
supplied.

\section{Independence of the non-zero periodic vectors}

The scope restriction of the section ``Scope restriction: independence without
spanning'' above concerned two separate weakenings: independence only along a
common-multiple subsequence, and the missing spanning clause. The
first is now removed: independence of the non-zero members of the family at
every sufficiently large length is
formalized.\footnote{Formal declaration:
\leanid{MPSTensor.periodicBasis_nonzero_subfamily_eventuallyLinearlyIndependent}
in \path{NonzeroSubfamily.lean}.} Only the spanning clause remains
unformalized.


\begin{thebibliography}{9}
\bibitem{DeLasCuevas2017Irreducible}
G.~de las Cuevas, J.~I.~Cirac, N.~Schuch, and D.~P\'erez-Garc\'ia,
\emph{Irreducible forms of matrix product states: Theory and applications},
J.~Math.~Phys.\ \textbf{58}, 121901 (2017), arXiv:1708.00029.

\bibitem{Cirac2016MPDO}
J.~I.~Cirac, D.~P\'erez-Garc\'ia, N.~Schuch, and F.~Verstraete,
\emph{Matrix product density operators: Renormalization fixed points and
boundary theories}, Ann.\ Phys.\ \textbf{378}, 100--149 (2017),
arXiv:1606.00608 (``Ci15'' in arXiv:1708.00029).

\bibitem{Wolf2012Quantum}
M.~M.~Wolf, \emph{Quantum Channels and Operations: Guided Tour}, unpublished
lecture notes (2012).
\end{thebibliography}

\end{document}
